
\section{Lineare Abbildungen}
In diesem Abschnitt werden wir eine spezielle Art von geometrischen Abbildungen untersuchen. Dazu werden wir Punkte in der Ebene oder im dreidimensionalen Raum durch ihre Ortsvektoren angeben. Der Punkt \( P = (2, 3) \) z.B. wird dann durch den Vektor \( \vec{p} = \begin{pmatrix} 2 \\ 3 \end{pmatrix} \) angegeben. So können wir Abbildungen durch ihre Wirkung auf die Ortsvektoren beschreiben. Sei z.B. \( T \) die Translation (=Parallelverschiebung) um den Vektor \( \begin{pmatrix} 3 \\ 1 \end{pmatrix} \). Unter dieser Verschiebung wird \( P = (2, 3) \) auf \( P' = (5, 4) \) abgebildet. Das schreiben wir so:

\[
T \begin{pmatrix} 2 \\ 3 \end{pmatrix} = \begin{pmatrix} 5 \\ 4 \end{pmatrix}
\]



\begin{center}
    \begin{tikzpicture}
        % Achsen
        \draw[->, thick] (-1,0) -- (6,0) node[below] {\small x};
        \draw[->, thick] (0,-1) -- (0,5) node[left] {\small y};
    
        % Punkte
        \filldraw[blue] (2,3) circle (1pt) node[left] {\small \(P(2,3)\)};
        \filldraw[red] (5,4) circle (1pt) node[right] {\small \(P'(5,4)\)};
        % Ortsvektoren
        \draw[->, dashed, blue, thick] (0,0) -- (2,3) node[midway, left] {\small \(\vec{p}\)};
        \draw[->, dashed, red, thick] (0,0) -- (5,4) node[midway, below right] {\small \(\vec{p'}\)};
    
        % Translationsvektor
        \draw[->] (2,3) -- (5,4) node[midway, above] {\small \(\begin{pmatrix} 3 \\ 1 \end{pmatrix}\)};
    
    \end{tikzpicture}
\end{center}

Häufig werden wir Abbildungen wie die Translation \( T \) durch ihre Gleichung angeben:

\[
T \begin{pmatrix} x \\ y \end{pmatrix} = \begin{pmatrix} x+3 \\ y+1 \end{pmatrix}
\]

\begin{LA}
{Abbildungsgleichungen}{Abbildungsgleichungen}
\begin{enumerate}
    \item Bestimmen Sie das Bild von \( \begin{pmatrix} 2 \\ 3 \end{pmatrix} \) und allgemein von \( \begin{pmatrix} x \\ y \end{pmatrix} \) unter der folgenden Abbildung:

\begin{enumerate}
    \item Verschiebung um \( \begin{pmatrix} -1 \\ 3 \end{pmatrix} \).
    % \[
    % T \begin{pmatrix} x \\ y \end{pmatrix} = \begin{pmatrix} x - 1 \\ y + 3 \end{pmatrix}
    % \]

    \item Reflexion (= Achsenspiegelung) an der $x$-Achse.
    % \[
    % R_x \begin{pmatrix} x \\ y \end{pmatrix} = \begin{pmatrix} x \\ -y \end{pmatrix}
    % \]

    \item Punktspiegelung am Ursprung \( O = (0,0) \).
    % \[
    % S_O \begin{pmatrix} x \\ y \end{pmatrix} = \begin{pmatrix} -x \\ -y \end{pmatrix}
    % \]

    \item Punktspiegelung an \( (1, -2) \).
    % \[
    % S_{(1,-2)} \begin{pmatrix} x \\ y \end{pmatrix} = \begin{pmatrix} 2 - x \\ -4 - y \end{pmatrix}
    % \]

    \item Rotation (= Drehung) um \( 90^\circ \) im Gegenuhrzeigersinn um den Ursprung.
    % \[
    % R_{90} \begin{pmatrix} x \\ y \end{pmatrix} = \begin{pmatrix} -y \\ x \end{pmatrix}
    % \]

    \item Streckung um den Faktor \( k = 3 \) mit Zentrum im Ursprung.
    % \[
    % S_3 \begin{pmatrix} x \\ y \end{pmatrix} = \begin{pmatrix} 3x \\ 3y \end{pmatrix}
    % \]

    \item Spiegelung an der Geraden \( x = 3 \).
    % \[
    % S_{x=3} \begin{pmatrix} x \\ y \end{pmatrix} = \begin{pmatrix} 6 - x \\ y \end{pmatrix}
    % \]
\end{enumerate}

\item Bestimmen Sie das Bild von \( \begin{pmatrix} 1 \\ 2 \\ 3 \end{pmatrix} \) und allgemein von \( \begin{pmatrix} x \\ y \\ z \end{pmatrix} \) unter den folgenden Abbildungen:

\begin{enumerate}
    \item Spiegelung am Ursprung \( O = (0,0,0) \).
    % \[
    % S_O \begin{pmatrix} x \\ y \\ z \end{pmatrix} = \begin{pmatrix} -x \\ -y \\ -z \end{pmatrix}
    % \]

    \item Spiegelung an der $yz$-Ebene.
    % \[
    % S_{yz} \begin{pmatrix} x \\ y \\ z \end{pmatrix} = \begin{pmatrix} -x \\ y \\ z \end{pmatrix}
    % \]

    \item Orthogonale (= senkrechte) Projektion auf die $xy$-Ebene.
    % \[
    % P_{xy} \begin{pmatrix} x \\ y \\ z \end{pmatrix} = \begin{pmatrix} x \\ y \\ 0 \end{pmatrix}
    % \]
\end{enumerate}
\end{enumerate}

\end{LA}

\clearpage
In diesem Kapitel werden wir uns mit einer speziellen Art von Abbildungen der Ebene oder des Raumes befassen. Unter „der Ebene“ bzw. „dem Raum“ verstehen wir dabei Folgendes:


\begin{definition}
    {$\R^2$ und $\R^3$}{}
    Die Menge aller Vektoren der Form \( \vec{v} = \begin{pmatrix} x \\ y \end{pmatrix} \), wobei \( x \) und \( y \) reelle Zahlen sind, nennt man \( \mathbb{R}^2 \):

    \[
    \mathbb{R}^2 = \left\{ \begin{pmatrix} x \\ y \end{pmatrix} \mid x \in \mathbb{R} \text{ und } y \in \mathbb{R} \right\}
    \]

    Analog gilt:

    \[
    \mathbb{R}^3 = \left\{ \begin{pmatrix} x \\ y \\ z \end{pmatrix} \mid x \in \mathbb{R} \text{ und } y \in \mathbb{R} \text{ und } z \in \mathbb{R} \right\}
    \]

\end{definition}
	
Unter allen Abbildungen der Ebene (oder des Raums) in sich gibt es einige mit einer speziellen Eigenschaft: die sogenannten linearen Abbildungen.

\begin{definition}
    {Lineare Abbildung}{}
    Eine Abbildung \( f \) heisst linear, wenn für alle Vektoren \( \vec{v} \) und \( \vec{w} \) gilt:

        \[
        f(\vec{v} + \vec{w}) = f(\vec{v}) + f(\vec{w})
        \]

        und für alle Zahlen \( \lambda \in \mathbb{R} \) gilt:

        \[
        f(\lambda \cdot \vec{v}) = \lambda \cdot f(\vec{v})
        \]

\end{definition}

\begin{bsp}
{}{}
Die Spiegelung an der $y$-Achse ist eine lineare Abbildung.

\vspace*{7cm}

% Die Spiegelung an der y-Achse in der Ebene hat die Gleichung:

% \[
% f: \mathbb{R}^2 \to \mathbb{R}^2, \quad f \begin{pmatrix} x \\ y \end{pmatrix} = \begin{pmatrix} -x \\ y \end{pmatrix}
% \]

% Die Spiegelung an der y-Achse ist eine lineare Abbildung, denn für alle Vektoren 
% \( \vec{v} = \begin{pmatrix} x_v \\ y_v \end{pmatrix} \) und 
% \( \vec{w} = \begin{pmatrix} x_w \\ y_w \end{pmatrix} \) gilt:

% \[
% f(\vec{v} + \vec{w}) = f \left( \begin{pmatrix} x_v \\ y_v \end{pmatrix} + \begin{pmatrix} x_w \\ y_w \end{pmatrix} \right)
% \]

% \[
% = f \begin{pmatrix} x_v + x_w \\ y_v + y_w \end{pmatrix} = \begin{pmatrix} -(x_v + x_w) \\ y_v + y_w \end{pmatrix}
% \]

% \[
% = \begin{pmatrix} -x_v - x_w \\ y_v + y_w \end{pmatrix} = \begin{pmatrix} -x_v \\ y_v \end{pmatrix} + \begin{pmatrix} -x_w \\ y_w \end{pmatrix}
% \]

% \[
% = f(\vec{v}) + f(\vec{w})
% \]

% und 

% \[
% f(\lambda \cdot \vec{v}) = f \left( \lambda \cdot \begin{pmatrix} x \\ y \end{pmatrix} \right) = f \begin{pmatrix} \lambda \cdot x \\ \lambda \cdot y \end{pmatrix}
% \]

% \[
% = \begin{pmatrix} -\lambda \cdot x \\ \lambda \cdot y \end{pmatrix} = \begin{pmatrix} \lambda \cdot (-x) \\ \lambda \cdot y \end{pmatrix} = \lambda \cdot \begin{pmatrix} -x \\ y \end{pmatrix}
% \]

% \[
% = \lambda \cdot f(\vec{v})
% \]

% Damit ist gezeigt, dass die Spiegelung an der y-Achse eine lineare Abbildung ist.

\end{bsp}
\clearpage

\begin{ueb}{Zu lineare Abbildungen}{lineareAbbildung}
    \begin{enumerate}
        \item Beschreiben Sie die folgende Abbildung in Worten (Spiegelung, Verschiebung, Drehung, Projektion, ...). Ist sie linear?
        \begin{enumerate}
            \item \( f \begin{pmatrix} x \\ y \end{pmatrix} = \begin{pmatrix} x+1 \\ y+1 \end{pmatrix} \)
            \item \( f \begin{pmatrix} x \\ y \end{pmatrix} = \begin{pmatrix} y \\ x \end{pmatrix} \)
            \item \( f \begin{pmatrix} x \\ y \end{pmatrix} = \begin{pmatrix} 2x \\ 2y \end{pmatrix} \)
            \item \( f \begin{pmatrix} x \\ y \end{pmatrix} = \begin{pmatrix} x \\ 0 \end{pmatrix} \)
            \item \( f \begin{pmatrix} x \\ y \\ z \end{pmatrix} = \begin{pmatrix} x \\ -y \\ z \end{pmatrix} \)
            \item \( f \begin{pmatrix} x \\ y \\ z \end{pmatrix} = \begin{pmatrix} 0 \\ y \\ z \end{pmatrix} \)
        \end{enumerate}
    
        \item Für lineare Abbildungen gilt für alle $\lambda \in \R$, dass \( f(\lambda \cdot \vec{v}) = \lambda \cdot f(\vec{v}) \) ist. Wenn man \( \lambda = 0 \) wählt, folgt, dass lineare Abbildungen den Nullvektor stets auf sich selbst abbilden. Nutzen Sie diese Eigenschaft, um drei Abbildungen in \Cref{LA:Abbildungsgleichungen} zu finden, die nicht linear sind.
    
        \item Sei \( f \) eine lineare Abbildung mit 
        \[
        f \begin{pmatrix} 1 \\ 0 \end{pmatrix} = \begin{pmatrix} 2 \\ 3 \end{pmatrix}, 
        \quad f \begin{pmatrix} 0 \\ 1 \end{pmatrix} = \begin{pmatrix} -1 \\ 1 \end{pmatrix}.
        \]
        Was ist dann \( f \begin{pmatrix} 4 \\ 5 \end{pmatrix} \)?  
        Tipp: Schreibe \( \begin{pmatrix} 4 \\ 5 \end{pmatrix} \) als  
        \[
        \begin{pmatrix} 4 \\ 5 \end{pmatrix} = 4 \cdot \begin{pmatrix} 1 \\ 0 \end{pmatrix} + 5 \cdot \begin{pmatrix} 0 \\ 1 \end{pmatrix}
        \]
        und verwende die Linearität.
    
      
            \item Sei \( f \) eine lineare Abbildung mit  
            \[
            f \begin{pmatrix} 2 \\ 1 \end{pmatrix} = \begin{pmatrix} -1 \\ 1 \end{pmatrix}, 
            \quad f \begin{pmatrix} 3 \\ -2 \end{pmatrix} = \begin{pmatrix} 2 \\ 4 \end{pmatrix}.
            \]
            Was ist dann \( f \begin{pmatrix} 1 \\ 11 \end{pmatrix} \)?  
            Tipp: Gehen Sie analog zu Aufgabe 3) vor.
    
            \item Challenge: Sei  
            \[
            f \begin{pmatrix} x \\ y \end{pmatrix} =
            \begin{cases} 
                \begin{pmatrix} 0 \\ x^2 / y \end{pmatrix}, & \text{falls } y \neq 0 \\[8pt]
                \vec{0}, & \text{falls } y = 0
            \end{cases}
            \]
            Zeigen Sie, dass \( f \) die Eigenschaft \( f(\lambda \cdot \vec{v}) = \lambda \cdot f(\vec{v}) \) erfüllt, aber nicht die Eigenschaft  
            \[
            f(\vec{v} + \vec{w}) = f(\vec{v}) + f(\vec{w}).
            \]
            Tipp: Beim Prüfen der Eigenschaft \( f(\lambda \cdot \vec{v}) = \lambda \cdot f(\vec{v}) \) müssen Sie drei Fälle unterscheiden:
            \begin{enumerate}
                \item \( \lambda = 0 \),
                \item \( \lambda \neq 0 \) und \( y = 0 \),
                \item \( \lambda \neq 0 \) und \( y \neq 0 \).
            \end{enumerate}
        
    \end{enumerate}
    
\end{ueb}
\clearpage
\emph{Lösungen}
\begin{enumerate}
    \item Abbildungen beschreiben...
    \begin{enumerate}
        \item Verschiebung um \( \begin{pmatrix} 1 \\ 1 \end{pmatrix} \), nicht linear.
        \item Spiegelung an \( y = x \), linear.
        \item Streckung um den Faktor \( k=2 \) mit Zentrum im Ursprung, linear.
        \item Orthogonale Projektion auf die $x$-Achse, linear.
        \item Spiegelung an der $xz$-Ebene, linear.
        \item Orthogonale Projektion auf die $yz$-Ebene, linear.
    \end{enumerate}
    
    \item Siehe oben.

    \item Berechnung:
    \[
    f \begin{pmatrix} 4 \\ 5 \end{pmatrix} = \begin{pmatrix} 3 \\ 17 \end{pmatrix}
    \]

    \item Lösung für \( f \begin{pmatrix} 1 \\ 11 \end{pmatrix} \):
    \[
    \begin{pmatrix} 1 \\ 11 \end{pmatrix} = a \cdot \begin{pmatrix} 2 \\ 1 \end{pmatrix} + b \cdot \begin{pmatrix} 3 \\ -2 \end{pmatrix} \Rightarrow a = 5, \quad b = -3
    \]

    \[
    f \begin{pmatrix} 1 \\ 11 \end{pmatrix} = 5 \cdot f \begin{pmatrix} 2 \\ 1 \end{pmatrix} - 3 \cdot f \begin{pmatrix} 3 \\ -2 \end{pmatrix} = \begin{pmatrix} -11 \\ -7 \end{pmatrix}
    \]

    \item Behauptung 1: \( f(\lambda \cdot \vec{v}) = \lambda \cdot f(\vec{v}) \)
    
    Beweis: Fallunterscheidung
    \begin{itemize}
        \item Falls \( \lambda = 0 \):
        \[
        f(\lambda \cdot \vec{v}) = f(\vec{0}) = \vec{0}, \quad \lambda f(\vec{v}) = 0 \cdot f(\vec{v}) = \vec{0}
        \]
        
        \item Falls \( \lambda \neq 0, y = 0 \):
        \[
        f(\lambda \cdot \vec{v}) = \vec{0}, \quad \lambda f(\vec{v}) = \vec{0}
        \]
        
        \item Falls \( \lambda \neq 0, y \neq 0 \):
        \[
        f(\lambda \cdot \vec{v}) = f \begin{pmatrix} \lambda x \\ \lambda y \end{pmatrix} = \begin{pmatrix} 0 \\ \frac{(\lambda x)^2}{\lambda y} \end{pmatrix} = \begin{pmatrix} 0 \\ \lambda \frac{x^2}{y} \end{pmatrix} = \lambda f(\vec{v})
        \]
    \end{itemize}

    Behauptung 2: \( f(\vec{v} + \vec{w}) = f(\vec{v}) + f(\vec{w}) \) stimmt nicht immer.

    Beweis: Es reicht ein Gegenbeispiel:
    \[
    f \left( \begin{pmatrix} 1 \\ 1 \end{pmatrix} + \begin{pmatrix} 0 \\ -1 \end{pmatrix} \right) = f \begin{pmatrix} 1 \\ 0 \end{pmatrix} = \vec{0} \neq f \begin{pmatrix} 1 \\ 1  \end{pmatrix} + f \begin{pmatrix} 0\\ -1  \end{pmatrix} = \begin{pmatrix} 0 \\ 1  \end{pmatrix} +  \begin{pmatrix} 0\\ 0  \end{pmatrix} = \begin{pmatrix}
        0\\1
    \end{pmatrix}
    \]
\end{enumerate}
